\chapter{2020年上海卷(春)}

\section{填空题}

\begin{problem}
集合 \(A=\{1,3\}\),\(B=\{1,2,a\}\),若 \(A\subseteq B\),则 \(a=\fillinblank{}\).
\end{problem}

\begin{problem}
不等式 \(\dfrac{1}{x}>3\) 的解集为 \(\fillinblank{}\).
\end{problem}

\begin{problem}
函数 \(y=\tan 2x\) 的最小正周期为 \(\fillinblank{}\).
\end{problem}

\begin{problem}
已知复数 \(z\) 满足 \(z+2\bar z=6+\i\),则 \(z\) 的实部为 \(\fillinblank{}\).
\end{problem}

\begin{problem}
已知 \(3\sin 2x=2\sin x\),\(x\in(0,\pi)\),则 \(x=\fillinblank{}\).
\end{problem}

\begin{problem}
若函数 \(y=a\cdot3^x+\dfrac{1}{3^x}\) 为偶函数,则 \(a=\fillinblank{}\).
\end{problem}

\begin{problem}
已知直线 \(l_1:x+ay=1\),\(l_2:ax+y=1\),若 \(l_1\parallel l_2\),则 \(l_1\) 与 \(l_2\) 的距离为 \(\fillinblank{}\).
\end{problem}

\begin{problem}
已知二项式 \((2x+\sqrt{x})^5\),则展开式中 \(x^3\) 的系数为 \(\fillinblank{}\).
\end{problem}

\begin{problem}
\(\triangle ABC\) 中,\(D\) 是 \(BC\) 中点,\(AB=2\),\(BC=3\),\(AC=4\),则 \(\overrightarrow{AD}\cdot\overrightarrow{AB}=\fillinblank{}\).
\end{problem}

\begin{problem}
已知 \(A=\{-3,-2,-1,0,1,2,3\}\),\(a\),\(b\in A\),则 \(|a|<|b|\) 的情况有 \(\fillinblank{}\) 种.
\end{problem}

\begin{problem}
已知 \(A_1\),\(A_2\),\(A_3\),\(A_4\),\(A_5\) 五个点,满足 \(\overrightarrow{A_nA_{n+1}}\cdot\overrightarrow{A_{n+1}A_{n+2}}=0\) (\(n=1\),\(2\),\(3\)),\(|A_nA_{n+1}|\cdot|A_{n+1}A_{n+2}|=n+1\) (\(n=1\),\(2\),\(3\)),则 \(|A_1A_5|\) 的最小值为 \(\fillinblank{}\).
\end{problem}

\begin{problem}
已知 \(f(x)=\sqrt{x-1}\),其反函数为 \(f^{-1}(x)\),若 \(f^{-1}(x)-a=f(x+a)\) 有实数根,则 \(a\) 的取值范围为 \(\fillinblank{}\).
\end{problem}

\section{单选题}

\begin{problem}
计算 \(\displaystyle\lim_{n\to\infty}\dfrac{3^n+5^n}{3^{n-1}+5^{n-1}}=\)(\quad).

\choices{\(3\)}{\(\dfrac{5}{3}\)}{\(\dfrac{3}{5}\)}{\(5\)}
\end{problem}

\begin{problem}
``\(\alpha=\beta\)''是``\(\sin^2\alpha+\cos^2\beta=1\)''的(\quad).

\choices{充分非必要条件}{必要非充分条件}{充要条件}{既非充分又非必要条件}
\end{problem}

\begin{problem}
已知椭圆 \(\dfrac{x^2}{2}+y^2=1\),作垂直于 \(x\) 轴的垂线交椭圆于 \(A\),\(B\) 两点,作垂直于 \(y\) 轴的垂线交椭圆于 \(C\),\(D\) 两点,且 \(AB=CD\),两垂线相交于点 \(P\),则点 \(P\) 的轨迹是(\quad).

\choices{椭圆}{双曲线}{圆}{抛物线}
\end{problem}

\begin{problem}
数列 \(\{a_n\}\) 各项均为实数,对任意 \(n\in\mathbb N^*\) 满足 \(a_{n+3}=a_n\),且行列式 \(\begin{vmatrix}a_n&a_{n+1}\\a_{n+2}&a_{n+3}\end{vmatrix}=c\) 为定值,则下列选项中不可能的是(\quad).

\choices{\(a_1=1,c=1\)}{\(a_1=2,c=2\)}{\(a_1=-1,c=4\)}{\(a_1=2,c=0\)}
\end{problem}

\section{解答题}

\begin{problem}
已知四棱锥 \(P-ABCD\),底面 \(ABCD\) 为正方形,边长为 \(3\),\(PD\perp\) 平面 \(ABCD\).
\begin{enumerate}
\item 若 \(PC=5\),求四棱锥 \(P-ABCD\) 的体积;
\item 若直线 \(AD\) 与 \(BP\) 的夹角为 \(60^\circ\),求 \(PD\) 的长.
\end{enumerate}

\FigureLayoutDeclare{img/2020/shanghai_spring/q17_fig1.png}{4.5cm}{13bp 0bp 25bp 0bp}
\bitmapfigure[max width=0.50\linewidth]{img/2020/shanghai_spring/q17_fig1.png}
\end{problem}

\begin{problem}
已知各项均为正数的数列 \(\{a_n\}\),其前 \(n\) 项和为 \(S_n\),\(a_1=1\).
\begin{enumerate}
\item 若数列 \(\{a_n\}\) 为等差数列,\(S_{10}=70\),求数列 \(\{a_n\}\) 的通项公式;
\item 若数列 \(\{a_n\}\) 为等比数列,\(a_4=\dfrac{1}{8}\),求满足 \(S_n>100a_n\) 时 \(n\) 的最小值.
\end{enumerate}
\end{problem}

\begin{problem}
有一条长为 \(120\) 米的步行道 \(OA\),\(A\) 是垃圾投放点 \(\omega_1\),若以 \(O\) 为原点,\(OA\) 为 \(x\) 轴正半轴建立直角坐标系,设点 \(B(x,0)\),现要建设另一座垃圾投放点 \(\omega_2(t,0)\),函数 \(f_t(x)\) 表示与 \(B\) 点距离最近的垃圾投放点的距离.
\begin{enumerate}
\item 若 \(t=60\),求 \(f_{60}(10)\),\(f_{60}(80)\),\(f_{60}(95)\) 的值,并写出 \(f_{60}(x)\) 的函数解析式;
\item 若可以通过 \(f_t(x)\) 与坐标轴围成的面积来测算扔垃圾的便利程度,面积越小越便利. 问:垃圾投放点 \(\omega_2\) 建在何处才能比建在中点时更加便利?
\end{enumerate}
\end{problem}

\begin{problem}
已知抛物线 \(y^2=x\) 上的动点 \(M(x_0,y_0)\),过 \(M\) 分别作两条直线交抛物线于 \(P\),\(Q\) 两点,交直线 \(x=t\) 于 \(A\),\(B\) 两点.
\begin{enumerate}
\item 若点 \(M\) 纵坐标为 \(\sqrt{2}\),求 \(M\) 与焦点的距离;
\item 若 \(t=-1\),\(P(1,1)\),\(Q(1,-1)\),求证:\(y_A\cdot y_B\) 为常数;
\item 是否存在 \(t\),使得 \(y_A\cdot y_B=1\) 且 \(y_P\cdot y_Q\) 为常数?若存在,求出 \(t\) 的所有可能值;若不存在,请说明理由.
\end{enumerate}
\end{problem}

\begin{problem}
已知非空集合 \(A\subseteq\R\),函数 \(y=f(x)\) 的定义域为 \(D\),若对任意 \(t\in A\) 且 \(x\in D\),不等式 \(f(x)\geq f(x+t)\) 恒成立,则称函数 \(f(x)\) 具有 \(A\) 性质.
\begin{enumerate}
\item 当 \(A=\{-1\}\),判断 \(f(x)=-x\),\(g(x)=2^x\) 是否具有 \(A\) 性质;
\item 当 \(A=(0,1)\),\(f(x)=x+\dfrac{1}{x}\),\(x\in[a,+\infty)\),若 \(f(x)\) 具有 \(A\) 性质,求 \(a\) 的取值范围;
\item 当 \(A=\{-2,m\}\),\(m\in\Z\),若 \(D\) 为整数集且具有 \(A\) 性质的函数均为常值函数,求所有符合条件的 \(m\) 的值.
\end{enumerate}
\end{problem}

